\section*{Chapitre \ref{ch:iv1:order-statistics} --- Statistiques d'ordre}
\markboth{\protect\omTocMark{}{Réponses~: Statistiques d'ordre}}{\protect\omTocMark{}{Réponses~: Statistiques d'ordre}}

\begin{solution}{iq:iv1:order-statistics:1}
Le plus grand a pour moyenne $\frac{n}{n + 1}$, le plus petit $\frac{1}{n + 1}$ et l'étendue $\frac{n - 1}{n + 1}$. Pour $n = 10$, le plus grand est inférieur à $0.9$ avec la probabilité $0.9^{10} \approx 0.349$.

Les $n$ points découpent l'intervalle en $n + 1$ morceaux, et ces morceaux sont échangeables\index{espacements uniformes}, de sorte que chacun a pour moyenne $\frac{1}{n + 1}$. Le plus petit est le premier morceau, le plus grand est tout sauf le dernier morceau, et l'étendue est tout sauf les deux morceaux extrêmes. Le plus grand est inférieur à une valeur exactement lorsque les $n$ points le sont tous.

Le cas $n = 1$ vérifie les trois résultats~: le plus grand et le plus petit ont tous deux pour moyenne $\tfrac12$, et l'étendue vaut 0.

\iqwrong{Donner $1 - \tfrac1n$ pour le plus grand. C'est le seuil en dessous duquel un tirage unique tombe avec la probabilité $1 - 1/n$, et non la moyenne du plus grand. L'écart est faible pour $n$ grand, mais l'erreur est flagrante pour $n = 1$.}
\end{solution}

\begin{solution}{iq:iv1:order-statistics:2}
La densité vaut $6x(1 - x)$, la moyenne $\tfrac12$ et la variance $1/20$. La moyenne est meilleure, avec une variance de $1/36$.

La médiane est proche de $x$ lorsqu'un point est proche de $x$, un en dessous et un au-dessus. Il y a $3!$ façons d'affecter ces rôles, ce qui donne la densité $6x(1 - x)$. En général, le $k$-ième plus petit de $n$ uniformes suit une loi Beta$(k, n - k + 1)$\index{loi Bêta des statistiques d'ordre}, ici Beta$(2, 2)$. Par symétrie, la moyenne vaut $\tfrac12$, et $E[M^2] = \int_0^1 6x^3(1 - x)\,dx = 3/10$, donc la variance vaut $3/10 - 1/4 = 1/20$.

La moyenne de trois variables a pour variance $\frac{1}{12 \cdot 3} = 1/36$, de sorte que la médiane est $1.8$ fois plus variable.

\iqwrong{Affirmer que la médiane est toujours le choix le plus robuste, donc le meilleur. La robustesse protège contre les valeurs aberrantes, et une loi uniforme n'en comporte aucune. Pour une distribution à bornes franches, les valeurs extrêmes sont au contraire les plus informatives.}
\end{solution}

\begin{solution}{iq:iv1:order-statistics:3}
Le plus rapide appartient au desk B avec la probabilité $2/3$, et les trois plus rapides appartiennent tous à ce desk avec la probabilité $\binom{20}{3}/\binom{30}{3} = 57/203 \approx 0.281$. La loi n'a pas d'importance.

Les 30 valeurs proviennent toutes d'une même loi continue, donc chaque ordre les concernant est aussi probable que les autres\index{echangeabilite des rangs@échangeabilité des rangs}. Les étiquettes A et B sont donc une affectation aléatoire de 10 et de 20 positions dans la liste triée. La probabilité que les trois premières positions reviennent toutes à B vaut la chaîne $\tfrac{20}{30} \cdot \tfrac{19}{29} \cdot \tfrac{18}{28}$.

C'est sur ce raisonnement que reposent les tests de rang. Si les trois plus rapides provenaient tous d'un même desk, on peut calculer cette probabilité sous l'hypothèse «~même loi~», sans aucun modèle de latence.

\iqwrong{Élever $2/3$ au cube, ce qui donne environ $0.296$. Cela traite les positions comme indépendantes, alors qu'elles sont tirées sans remise.}
\iqfollow{Quel est le test auquel cela conduit~?}{Le test des rangs de Mann--Whitney~: il additionne les rangs d'un groupe et les compare à la distribution de toutes les affectations possibles.}
\end{solution}

\begin{solution}{iq:iv1:order-statistics:4}
L'estimation sans biais vaut $60 \cdot (1 + \tfrac15) - 1 = 71$. Le double de la moyenne moins un a une variance $7/3$ fois plus grande, soit environ $18^2$ contre $12^2$ pour $N = 71$.

Le plus grand de $k$ numéros de série a $P(\max = x) = \binom{x - 1}{k - 1}/\binom{N}{k}$, et la sommation donne $E[\max] = k(N + 1)/(k + 1)$\index{problème tank allemand}. Donc $N = \frac{k + 1}{k} E[\max] - 1$, et en remplaçant $E[\max]$ par ce que l'on a observé, on obtient un estimateur sans biais.

Une image sous-tend ce résultat. Les $k$ numéros découpent $1, \dots, N$ en $k + 1$ intervalles de même longueur moyenne. Les numéros au-dessus du maximum forment un intervalle, et la moyenne des $k$ intervalles visibles vaut environ $m/k$. On ajoute donc cette quantité à $m$.

La moyenne empirique a pour moyenne $(N + 1)/2$~: doubler cette moyenne et retrancher un donne donc aussi un estimateur sans biais, mais le numéro le plus élevé porte à lui seul toute l'information de l'échantillon. Connaissant le maximum, les autres numéros sont simplement uniformes en dessous de celui-ci et n'apportent plus rien sur $N$\index{statistique exhaustive}. La règle de la moyenne l'ignore, et avec un échantillon $1, 2, 3, 4, 60$ elle donne 27, ce qui est inférieur à un numéro de série effectivement observé.

\iqwhy{Savoir quel numéro unique d'un échantillon porte l'information, et ne pas la diluer par une moyenne. Le même raisonnement s'applique dès qu'une série possède une borne franche, comme une limite de position ou un plancher de prix.}
\end{solution}

\begin{solution}{iq:iv1:order-statistics:5}
Les variances valent $\theta^2/30$ et $\theta^2/120$ pour $n = 10$, de sorte que la moyenne est 4 fois plus variable. Le meilleur multiple est $\frac{n + 2}{n + 1}\max X_i$, avec une erreur quadratique moyenne $\theta^2/(n + 1)^2$, soit ici $\theta^2/121$ contre $\theta^2/120$.

Prenons $\theta = 1$. La moyenne a pour variance $\frac{1}{12n}$, donc $2\bar{X}$ a pour variance $\frac{1}{3n}$. Le maximum vérifie $E[M] = \frac{n}{n + 1}$ et $E[M^2] = \frac{n}{n + 2}$, donc $\mathrm{Var}(M) = \frac{n}{(n + 1)^2(n + 2)}$, et le maximum mis à l'échelle a pour variance $\frac{1}{n(n + 2)}$. Le rapport vaut $(n + 2)/3$.

Remarquons les vitesses de convergence~: $1/n$ pour la moyenne contre $1/n^2$ pour le maximum. Une extrémité d'une distribution s'estime bien plus vite que son centre\index{superefficacité à une frontière}.

Pour $cM$, l'erreur quadratique moyenne vaut $c^2 \frac{n}{n + 2} - 2c \frac{n}{n + 1} + 1$, une parabole en $c$ dont le minimum se situe en $c = \frac{n + 2}{n + 1}$. On accepte ainsi un léger biais pour une variance plus faible, mais le gain est minime, $1/121$ contre $1/120$\index{compromis biais--variance}. Le maximum brut, l'estimation du maximum de vraisemblance, a une erreur de $\frac{2}{(n + 1)(n + 2)} = 1/66$, soit près du double.

\iqwrong{Affirmer que l'estimation du maximum de vraisemblance est la meilleure parce qu'elle l'est souvent. Ici elle est biaisée vers le bas par construction, puisqu'elle ne peut jamais dépasser le $\theta$ réel.}
\end{solution}

\begin{solution}{iq:iv1:order-statistics:6}
Le deuxième plus élevé de cinq ordres d'achat a pour moyenne $2/3$. Sachant que le plus élevé vaut $0.8$, il a pour moyenne $0.64$. Le deuxième plus élevé de quatre dés a pour moyenne $5313/1296 \approx 4.10$.

À partir des espacements\index{espacements uniformes}~: le deuxième plus élevé est tout sauf les deux plus hauts parmi les six morceaux de même moyenne, donc sa moyenne vaut $4/6$.

Sachant que le plus élevé vaut $h$, les quatre autres sont des uniformes indépendantes sur $[0, h]$, et le deuxième plus élevé est le plus grand d'entre elles, de moyenne $\tfrac45 h$. Pour $h = 0.8$, cela donne $0.64$. Pour vérifier les deux réponses, le plus élevé a pour densité $5h^4$, et en moyennant $\tfrac45 h$ selon cette densité on obtient $\tfrac45 \cdot \tfrac56 = \tfrac23$.

Pour les dés, utilisons la formule de la queue\index{formule de la somme des queues}. Le deuxième plus élevé est au moins égal à $k$ lorsqu'au moins deux des quatre dés le sont, donc avec $p = (7 - k)/6$,
\[
P(\text{deuxième} \ge k) = 1 - (1 - p)^4 - 4p(1 - p)^3,
\]
et la somme de ces probabilités pour $k = 1, \dots, 6$ donne $5313/1296$.

\iqwrong{Appliquer le résultat uniforme aux dés, en prenant $1 + 5 \cdot \tfrac35 = 4$. Un dé n'est pas une uniforme à échelle modifiée, et les ex æquo rendent la réponse légèrement différente.}
\iqwhy{Dans une enchère au second prix, le gagnant paie la deuxième plus haute valeur, donc cette moyenne est le revenu du vendeur. Connaître la plus haute enchère modifie ce calcul, car les autres sont alors bornées par elle.}
\end{solution}

\begin{solution}{iq:iv1:order-statistics:7}
Pour un échantillon gaussien, la variance de la médiane vaut $\pi/2 \approx 1.57$ fois celle de la moyenne. Pour la loi de Laplace, elle vaut la moitié de celle de la moyenne. Ainsi, avec des queues suffisamment épaisses, la médiane l'emporte.

La médiane empirique a pour variance asymptotique $\frac{1}{4 n f(m)^2}$, où $f(m)$ désigne la densité à la médiane vraie\index{variance asymptotique de la médiane}. La raison en est que, autour de la médiane, le nombre de points situés sous un niveau suit une loi binomiale de variance $n/4$, et que la pente du nombre moyen est $n f(m)$~; en divisant, on obtient l'erreur sur la position.

Loi normale centrée réduite~: $f(0) = 1/\sqrt{2\pi}$, donc la variance vaut $\frac{\pi}{2n}$, contre $\frac1n$ pour la moyenne.

Loi de Laplace~: $f(0) = \tfrac12$, donc la variance vaut $\frac1n$, alors que la moyenne a pour variance $\frac2n$.

\iqwrong{Affirmer que la médiane est toujours moins efficace parce qu'elle «~jette des données~». Cela dépend de la densité au centre comparée à celle des queues. Un pic central et des queues épaisses favorisent la médiane.}
\iqfollow{Et à l'extrême, avec des queues très épaisses~?}{Pour un échantillon de Cauchy, la moyenne ne vaut pas mieux qu'un tirage unique, quelle que soit la taille $n$, alors que la médiane a pour variance $\frac{\pi^2}{4n}$. Les rendements réels se situent entre la loi normale et ces cas extrêmes, ce qui explique qu'il vaille la peine de calculer une moyenne tronquée ou une médiane d'une série quotidienne en complément de la moyenne.}
\end{solution}

\begin{solution}{iq:iv1:order-statistics:8}
Le premier tombe en panne au bout de $1/5$ d'année en moyenne, le troisième après $47/60$, le dernier après $137/60$. Le moment de la première panne et le délai entre celle-ci et la dernière sont indépendants.

Cinq serveurs en fonctionnement tombent en panne à un taux total de 5, donc la première panne survient au bout de $1/5$ en moyenne. Par absence de mémoire\index{statistiques d'ordre exponentielles}, les quatre restants repartent alors à zéro, avec un taux de 4, et ainsi de suite. Les instants triés sont donc des sommes d'intervalles indépendants de moyennes $\tfrac15, \tfrac14, \tfrac13, \tfrac12, 1$.

Le troisième instant vaut $\tfrac15 + \tfrac14 + \tfrac13$, et le dernier est la somme des cinq intervalles. Le délai entre la première et la dernière panne est une somme des quatre autres intervalles, donc il est indépendant de la première, avec pour moyenne $25/12$.

Cette indépendance est une propriété particulière des lois exponentielles. Dans un échantillon uniforme, un minimum élevé laisse peu de place, de sorte que le minimum et l'étendue sont fortement corrélés négativement. Une simulation donne une corrélation proche de zéro pour les lois exponentielles, et inférieure à $-0.3$ pour les lois uniformes.

\iqwrong{Prendre un an pour le dernier, c'est-à-dire la moyenne d'une seule durée de vie. Le dernier de cinq doit survivre aux autres, et cela prend plus de deux fois plus de temps.}
\iqwhy{C'est le schéma classique des pannes matérielles avec redondance. Avec trois répliques, le temps jusqu'à la perte du dernier exemplaire ne croît que lentement avec la redondance, comme une somme harmonique, à moins que l'on ne répare entre-temps.}
\end{solution}

\begin{solution}{iq:iv1:order-statistics:9}
$P(N > n) = \frac{1}{n + 1}$. Le plus petit $n$ tel que $P(N \le n) \ge \frac12$ vaut 1, et l'espérance est infinie. Les événements de record sont indépendants.

$N > n$ signifie que le premier tirage est le plus grand parmi les $n + 1$ premiers. Toutes les permutations sont équiprobables, donc chacun des $n + 1$ tirages a la même probabilité d'être le plus grand\index{echangeabilite des rangs@échangeabilité des rangs}.

$P(N \le 1) = \tfrac12$ exactement, et $P(N \le 0) = 0$, donc le plus petit tel $n$ vaut 1. La médiane n'est pas unique, car $P(N \ge 2) = \tfrac12$ également~: tout nombre entre 1 et 2 est une médiane. L'espérance est la somme des probabilités de queue, $\sum_{n \ge 0} \frac{1}{n + 1}$, qui diverge\index{formule de la somme des queues}. Si le premier tirage tombe par hasard très haut, l'attente est très longue, et ces cas rares dominent la moyenne.

Pour les records~: le fait que le tirage $k$ soit un record est une affirmation sur la position du plus grand des $k$ premiers. L'ordre relatif des 9 premiers tirages, en particulier la place du plus grand des 5 premiers, est indépendant du fait que le tirage 10 les dépasse tous\index{indépendance des événements de record}. La probabilité des deux vaut donc $\tfrac15 \cdot \tfrac{1}{10} = \tfrac{1}{50}$, ce qu'une simulation confirme.

\iqwrong{Dire que l'attente moyenne doit être finie parce que l'on termine presque toujours après quelques tirages. La probabilité d'une longue attente décroît trop lentement, comme $1/n$, pour que la somme converge.}
\iqwhy{Dans une suite de tirages indépendants, un nouveau plus haut historique devient de plus en plus rare. Une longue série sans nouveau plus haut est attendue sous ce modèle, et elle ne dit rien à elle seule d'un changement.}
\end{solution}

\begin{solution}{iq:iv1:order-statistics:10}
Entre le plus petit et le plus grand~: $1 - 2 \cdot 2^{-10} = 511/512 \approx 0.998$. Entre le troisième plus petit et le troisième plus grand~: $912/1024 \approx 0.891$.

Chaque tirage est sous la vraie médiane avec probabilité exactement $\tfrac12$, quelle que soit la loi\index{intervalle indépendant de la loi}. Le nombre $B$ de tirages sous la médiane suit donc une loi binomiale de paramètres $n = 10$ et $p = \tfrac12$.

La médiane dépasse le plus petit tirage sauf si $B = 0$, et elle est sous le plus grand sauf si $B = 10$, ce qui donne la première réponse. La médiane se situe entre le troisième plus petit et le troisième plus grand exactement lorsque au moins 3 tirages sont en dessous d'elle et au moins 3 au-dessus, c'est-à-dire $3 \le B \le 7$~:
\[
\frac{\binom{10}{3} + \binom{10}{4} + \binom{10}{5} + \binom{10}{6} + \binom{10}{7}}{2^{10}} = \frac{120 + 210 + 252 + 210 + 120}{1024} = \frac{912}{1024}.
\]
Une simulation avec des tirages exponentiels et normaux donne les mêmes deux nombres.

\iqwrong{Utiliser un intervalle normal sur l'échantillon, $\bar{X} \pm 2s/\sqrt{n}$. Cela suppose des hypothèses que l'on n'a pas. Pour une médiane, le comptage binomial donne un intervalle exact.}
\iqfollow{Cela fonctionne-t-il pour la moyenne~?}{Non. Une loi peut avoir la même médiane avec n'importe quelle moyenne, à cause d'une petite probabilité très éloignée. Aucun intervalle construit à partir d'un échantillon fini ne peut couvrir la moyenne sous toutes les lois.}
\end{solution}

\begin{solution}{iq:iv1:order-statistics:11}
Le meilleur ratio de Sharpe mesuré vaut environ $1.08$, puisque le plus grand de 1000 normales centrées réduites a pour espérance environ $3.24$. Le meilleur attendu croît comme $\sqrt{2 \ln n}$, très lentement.

Estimation à la main~: le plus grand de $n$ se situe au niveau dépassé avec probabilité $1/n$\index{valeur extrême d'échantillons gaussiens}. Ici, cela donne $1/1000$, soit $3.09$ écarts-types, et en prenant la médiane du maximum, de queue $\ln 2/n$, on obtient $3.2$. On multiplie par $\tfrac13$ et l'on obtient environ $1.07$. La moyenne exacte, issue de $\int x\, n\phi(x)\Phi(x)^{n - 1}dx$, vaut $3.24$, ce qu'une simulation confirme.

La formule $\sqrt{2 \ln 1000} \approx 3.72$ a la bonne croissance, mais elle surestime d'un demi-écart-type à cette taille. Passer de 100 à 1000 stratégies, ou de 1000 à 10\,000, ajoute seulement environ $0.6$ à $0.7$ écart-type.

Avec une seule stratégie, il n'y a aucune sélection, et le meilleur attendu vaut 0.

\iqwrong{Dire que 1.08 sur neuf ans est trop élevé pour être dû au hasard, car sa statistique $t$ dépasse 3. Ce test vaut pour une seule stratégie choisie à l'avance. Après 1000 essais, on doit s'attendre à cette valeur lorsqu'aucune stratégie n'a la moindre compétence.}
\iqwhy{C'est la première question à poser à un backtest~: combien de choses ont été essayées~? Un seuil corrigé exige le nombre d'essais, et le vrai nombre inclut les tentatives que personne n'a notées.}
\end{solution}

\begin{solution}{iq:iv1:order-statistics:12}
$P(R < r) = n r^{n - 1} - (n - 1) r^n$, et il faut $n = 46$ tirages, pour lesquels la probabilité de rester sous le seuil vaut $0.048$ (contre $0.052$ pour $n = 45$). Pour toute loi continue, avec une confiance de 95~\%, au moins 90~\% de la loi se situe entre le plus petit et le plus grand de 46 tirages.

Considérons l'étendue inférieure à $r$. Supposons que le plus petit point soit celui en $x$, choisi de $n$ façons, et que les $n - 1$ autres doivent tomber dans $[x, x + r]$. Pour $x \le 1 - r$, cela a pour probabilité $r^{n - 1}$, et pour $x$ au-delà de cette valeur, la probabilité vaut $(1 - x)^{n - 1}$. Donc
\[
P(R < r) = n\int_0^{1 - r} r^{n - 1}dx + n\int_{1 - r}^1 (1 - x)^{n - 1} dx = n r^{n - 1}(1 - r) + r^n,
\]
C'est bien la formule. Une vérification~: pour $n = 2$, elle donne $2r - r^2$, ce qui correspond à la loi de $|X - Y|$.

Pour une loi générale de fonction de répartition $F$, les valeurs $F(X_i)$ sont uniformes\index{transformation intégrale de probabilité}. La masse entre le plus petit et le plus grand échantillon est exactement l'étendue uniforme $F(\max) - F(\min)$, si bien que le résultat tient sans aucune hypothèse\index{intervalle de tolérance}.

\iqwrong{Résoudre $0.9^n = 0.05$, ce qui donne $n = 29$. C'est juste pour une couverture unilatérale, avec le seul maximum. Couvrir les deux extrémités demande davantage de tirages.}
\iqwhy{Cela fournit une bande de stress sans modèle~: 46 jours observés bornent 90~\% de la loi avec une confiance de 95~\%, quelle que soit l'allure des rendements, pourvu que les jours soient indépendants et tirés d'une même loi. Ce n'est souvent ni l'un ni l'autre.}
\end{solution}

\begin{solution}{iq:iv1:order-statistics:13}
Les variances valent $1/132$ pour la moyenne, $1/52$ pour la médiane et $1/312$ pour le milieu de l'étendue. Le milieu de l'étendue l'emporte, et sa variance décroît comme $1/n^2$.

Pour la moyenne, la variance vaut $\frac{1}{12n}$. La médiane de $n = 2k + 1$ uniformes suit une loi Beta$(k + 1, k + 1)$, de variance $\frac{1}{4(n + 2)}$, soit ici $1/52$. Pour le milieu de l'étendue, le maximum et le minimum ont la même variance $\frac{n}{(n + 1)^2(n + 2)}$ et une covariance $\frac{1}{(n + 1)^2(n + 2)}$, donc
\[
\mathrm{Var}\Bigl(\frac{\max + \min}{2}\Bigr) = \frac{1}{2} \cdot \frac{n + 1}{(n + 1)^2(n + 2)} = \frac{1}{2(n + 1)(n + 2)}.
\]
Une simulation confirme les trois à 1~\% près. Pour $n = 1$, les trois coïncident, à $1/12$, ce qui est une autre vérification.

Le milieu de l'étendue l'emporte parce que la densité présente des bords francs et reste plate à l'intérieur. La vraisemblance ne dépend que des deux extrêmes~; connaissant ceux-ci, les autres points n'apportent donc plus d'information sur $\theta$\index{statistique exhaustive}. La médiane est la moins bonne, car elle n'utilise que de l'information locale, en un endroit sans particularité.

\iqwrong{Recommander le milieu de l'étendue dans tous les cas. C'est le pire choix avec des queues épaisses, et sur un échantillon gaussien sa variance dépasse déjà celle de la moyenne. Une simulation à $n = 11$ la montre plus de $1.5$ fois plus élevée. Un estimateur convient à une forme, et il n'en existe aucun universellement bon.}
\end{solution}

\begin{solution}{iq:iv1:order-statistics:14}
Aucun couple de points n'est à moins de 10 ms, avec probabilité $0.91^{10} \approx 0.389$. L'écart minimal a pour espérance $\frac{1}{n^2 - 1} = 1/99$ de seconde, soit environ $10.1$ ms.

La voie directe est un argument de contraction\index{argument de contraction}. Prenons les points dans l'ordre croissant $x_1 < \dots < x_n$, tous avec des écarts d'au moins $d$. Posons $y_i = x_i - (i - 1)d$. C'est une bijection de tels ensembles vers les ensembles triés de $n$ points de $[0, 1 - (n - 1)d]$, qui conserve le volume. Donc la probabilité que tous les écarts soient d'au moins $d$ vaut
\[
\bigl(1 - (n - 1)d\bigr)^n,
\]
ici $(1 - 9 \cdot 0.01)^{10}$. La formule vaut 1 en $d = 0$ et 0 lorsque $d = 1/(n - 1)$, quand les points ne peuvent plus se loger.

L'espérance de l'écart minimal est l'intégrale de cette queue, $\int_0^{1/(n - 1)} (1 - (n - 1)t)^n dt = \frac{1}{(n - 1)(n + 1)}$. Ainsi, une probabilité d'environ $0.39$ à 10 ms côtoie une espérance d'environ 10 ms, et c'est une loi à longue queue à droite.

\iqwrong{Donner $(1 - 0.02)^{45} \approx 0.40$, en traitant les 45 paires comme indépendantes, chacune de probabilité $2d$ d'être proches. C'est une approximation pour les événements rares~; elle tombe près de la valeur exacte, mais ici les paires ne sont pas indépendantes, et elle échouerait pour un $d$ plus grand.}
\iqwhy{Les rafales surviennent au hasard. Dix ordres dans une seconde ont le plus souvent deux d'entre eux à moins de 10 ms l'un de l'autre, si bien qu'un moteur d'appariement (matching engine) ou un contrôle de risque doit être conçu pour cela, même si l'écart moyen vaut 100 ms.}
\end{solution}

\enlargethispage{-2\baselineskip}%
\begin{solution}{iq:iv1:order-statistics:15}
Le plus grand de trois a pour espérance $\frac{3}{2\sqrt{\pi}} \approx 0.846$, et le plus grand de quatre a pour espérance $\frac{3}{2\sqrt{\pi}}\bigl(1 + \frac{2}{\pi} \arcsin\frac13\bigr) \approx 1.029$. La méthode va jusqu'à cinq (environ $1.163$), et six exigerait la probabilité que quatre normales corrélées soient toutes positives, pour laquelle il n'existe pas de formule simple.

Partons de $E_n = n\int x\phi(x)\Phi(x)^{n - 1}dx$. Comme $x\phi = -\phi'$, on intègre par parties\index{intégration par parties de type Stein}~:
\[
E_n = n(n - 1) \int \phi(x)^2\, \Phi(x)^{n - 2} dx.
\]
Or, $\phi(x)^2 = \frac{1}{2\sqrt{\pi}}$ fois la densité de $W \sim N(0, \tfrac12)$. Donc
\[
E_n = \frac{n(n - 1)}{2\sqrt{\pi}}\, E\bigl[\Phi(W)^{n - 2}\bigr],
\]
et $\Phi(W)^{n - 2}$ est la probabilité que $n - 2$ nouvelles normales centrées réduites indépendantes $A_i$ soient toutes inférieures à $W$. Les différences $A_i - W$ ont pour variance $\tfrac32$, et deux quelconques d'entre elles ont pour covariance $\tfrac12$, donc une corrélation de $\tfrac13$\index{probabilité d'orthant}. La probabilité qu'elles soient toutes négatives est une probabilité d'orthant.

Pour $n = 3$~: une variable, de probabilité $\tfrac12$, donc $E_3 = \frac{6}{2\sqrt{\pi}} \cdot \tfrac12$.

Pour $n = 4$~: deux variables, de probabilité $\tfrac14 + \frac{1}{2\pi}\arcsin\frac13$. En multipliant par $\frac{12}{2\sqrt{\pi}}$, on obtient la formule annoncée.

Pour $n = 5$~: trois variables, de probabilité $\tfrac18 + \frac{3}{4\pi}\arcsin\frac13$, donc $E_5 = \frac{5}{4\sqrt{\pi}}\bigl(1 + \frac{6}{\pi}\arcsin\frac13\bigr)$. Ces trois résultats coïncident avec une intégrale numérique à neuf décimales.

Une seconde méthode pour $n = 3$ n'utilise aucune intégrale. Les trois distances deux à deux ont pour somme le double de l'étendue. Chacune a pour espérance $2/\sqrt{\pi}$, donc l'étendue a pour espérance $3/\sqrt{\pi}$, et le plus grand en a la moitié, car le plus petit a la même espérance au signe près.

\iqwrong{Utiliser $\sqrt{2 \ln n}$, ce qui donne $1.48$ pour $n = 3$. C'est une vitesse de croissance asymptotique, très éloignée de la vraie valeur pour les petites valeurs de $n$. Un quant doit connaître les petites valeurs~: $0.56$, $0.85$, $1.03$, $1.16$ pour $n = 2$ à $5$.}
\end{solution}
